\documentclass[11pt,a4paper]{article}
\usepackage[T2A]{fontenc}
\usepackage[utf8]{inputenc}
\usepackage[russian,english]{babel}
\usepackage{amsmath,amssymb,amsthm,mathtools}
\usepackage{setspace}
\usepackage{enumitem}
\usepackage[hidelinks]{hyperref}
\setstretch{1.3}
\frenchspacing
\hyphenpenalty=10000

\newcommand{\head}[1]{\par\noindent\textbf{\underline{#1}}\quad}
\newcommand{\sheettitle}[4]{% title, subtitle, homework-flag, date
  \begin{center}
  {\huge\bfseries #1\par}\vspace{\dimexpr1.2em-2.5pt\relax}
  {\Large #2\par}\vspace{\dimexpr1em-0.4pt\relax}
  \if\relax\detokenize{#3}\relax\else{\Large #3\par}\vspace{\dimexpr1.15em-2pt\relax}\fi
  {\foreignlanguage{russian}{#4}\par}
  \end{center}\vspace{-1.5pt}}
\newcommand{\src}[1]{\unskip\nobreak\hfil\penalty50\hskip1em\hbox{}\nobreak\hfill{\small(#1)}{\parfillskip=0pt\par}}
\begin{document}
\sheettitle{Combinatorics -- 4}{Graphs}{Homework}{14 сентября 2026 года}

\head{Theoretical minimum.} $\sum_v\deg v=2|E|$, so the number of odd-degree vertices is even. A
graph on $n$ vertices is a tree if and only if it is connected with $n-1$ edges, if and only if
it is acyclic with $n-1$ edges, if and only if any two vertices are joined by a unique path; there
are $n^{n-2}$ labelled trees (Cayley, by the Pr\"ufer code). A graph is bipartite exactly when it
has no odd cycle; a connected graph has an Eulerian circuit exactly when all degrees are even.
Hall's theorem: a bipartite graph has a matching saturating $A$ if and only if
$|N(S)|\geqslant|S|$ for all $S\subseteq A$. For a connected plane graph $v-e+f=2$, whence
$e\leqslant3v-6$, the non-planarity of $K_5$ and $K_{3,3}$, a vertex of degree at most $5$, and
$6$-colourability. The number of walks of length $k$ is read off from powers of the adjacency
matrix, whose largest eigenvalue lies between the average and the maximum degree and whose
spectrum is symmetric exactly for bipartite graphs (Linear algebra~-- 11).

\par\vspace{-1pt}\noindent\rule{\linewidth}{0.4pt}\par\vspace{5pt}

\begin{enumerate}
\item There are $n$ people in a city and each of them has exactly $1000$ friends, friendship being
symmetric. Prove that one can select a group $S$ of people such that at least $n/2017$ of the
members of $S$ have exactly two friends inside $S$. \src{IMC}

\item Let $T$ be a tree with $n$ vertices and let $d(u,v)$ be the distance between $u$ and $v$.
Put $W(T)=\sum_{\{u,v\}}d(u,v)$ and $H(T)=\sum_{\{u,v\}}\frac1{d(u,v)}$, the sums being over
unordered pairs of distinct vertices. Prove that $W(T)\,H(T)\geqslant\frac{(n-1)^3(n+2)}4$.
\src{IMC}

\item An alien race has three genders: male, female and emale. A married triple consists of three
persons, one of each gender, who all like each other; every person belongs to at most one triple,
and liking is mutual. An expedition has $n$ members of each gender, and every member likes at
least $k$ persons of each of the two other genders. (a) Show that if $n$ is even and $k=n/2$, it
may be impossible to form even one married triple. (b) Show that if $k\geqslant3n/4$, then a
perfect system of $n$ married triples always exists. \src{IMC}

\item A round-robin tournament among $2n$ teams lasted $2n-1$ days: on each day every team played
one game against another team, one winning and one losing, and over the tournament each team
played every other team exactly once. Can one necessarily choose one winning team from each day
without choosing any team more than once? \src{Putnam}

\item In a tournament $2011$ players meet $2011$ times; every game is played by all of them at
once and each player either wins or loses it. Two $2011\times2011$ matrices $T$ and $W$ start at
zero; after every game, for every pair $(h,k)$, including $h=k$, the entry $T_{hk}$ grows by $1$
if players $h$ and $k$ tied (both won or both lost), while if $h$ won and $k$ lost then $W_{hk}$
grows by $1$ and $W_{kh}$ drops by $1$. Prove that at the end $\det(T+iW)$ is a non-negative
integer divisible by $2^{2010}$. \src{Putnam 2011}

\item Let $G$ be a simple graph on $2n$ vertices such that for every partition of its vertex set
into two sets $V_1$ and $V_2$ of size $n$ the number of edges inside $V_1$ equals the number of
edges inside $V_2$. Prove that all vertices have equal degrees. \src{CIIM 2014}

\item Let $G$ be a connected graph on vertices $1,\dots,n$ with $n\geqslant3$. For each pair of
edges $ij$ and $jk$ sharing a vertex take the three-cycles $(ijk)$ and $(ikj)$ as permutations.
Prove that these cycles generate the whole alternating group $A_n$.
\src{\foreignlanguage{russian}{Мехмат} 2011}

\item Let $G$ be a connected simple graph that is not regular, on $n$ vertices with degrees $d_v$.
Prove that
\[
\sum_{\{u,v\}\subseteq V}(d_u-d_v)^2\geqslant
\begin{cases}n-1,&n\text{ odd},\\[4pt]2n-4,&n\text{ even},\end{cases}
\]
and determine when equality holds. \src{IMC 2025, shortlist}

\item Let $G$ be a graph on $n$ vertices with $e$ edges and degrees $d_v$. Prove that
\[
e\biggl(e+\binom{n-1}2-(n-1)\biggr)\geqslant(n-1)\sum_v\binom{d_v}2,
\]
and find all graphs attaining equality. (The spectral facts --- $\bar d\leqslant\lambda_1\leqslant
\Delta$, the symmetry of the spectrum for bipartite graphs, closed walks as traces --- are Linear
algebra~-- 11, homework 9.) \src{IMC 2025, shortlist}

\item Prove that in any graph the vertices can be split into two parts so that at least half the
edges go between them, and that a graph with $m$ edges has a bipartite subgraph with at least
$\frac m2+\sqrt{\frac m8+\frac1{64}}-\frac18$ edges. Compare with the probabilistic argument of
Probability~-- 4.
\end{enumerate}
\end{document}
